\section*{Chapter \ref{ch:b3:supramolecular} --- Supramolecular Chemistry}

\begin{solution}{exo:b3:supramolecular:1}
Ion--dipole; $\pi$ stacking; cation--$\pi$; \omterm{def:b3:supramolecular:interactions}{halogen bond} (I to N); three hydrogen bonds.
\end{solution}

\begin{solution}{exo:b3:supramolecular:2}
$\Delta G^\circ = -RT\ln K = -8.314 \times 298 \times \ln(\num{1.0e4}) = \qty{-22.8}{kJ/mol}$.
\end{solution}

\begin{solution}{exo:b3:supramolecular:3}
$K_{\ce{K+}}/K_{\ce{Na+}} = 10^{1.7} = 50$; the difference is $RT\ln 50 = \qty{9.7}{kJ/mol}$ in favour of potassium.
\end{solution}

\begin{solution}{exo:b3:supramolecular:4}
$n/(n + m) = 1/3$: one \omterm{def:b3:supramolecular:supramolecular}{host} for two \omterm{def:b3:supramolecular:supramolecular}{guests}, $\mathrm{HG_2}$.
\end{solution}

\begin{solution}{exo:b3:supramolecular:5}
$H_0 + G_0 + 1/K = \qty{4.667e-3}{mol/L}$; $[\mathrm{HG}] = (4.667 - \sqrt{4.667^2 - 4 \times 2.0 \times 2.0})/2 \times 10^{-3} = \qty{1.13e-3}{mol/L}$: bound
fraction 0.566, $\delta = 3.560 + 0.160 \times 0.566 = \qty{3.651}{ppm}$.
\end{solution}

\begin{solution}{exo:b3:supramolecular:6}
$\log K = 18.3 - 8.6 = 9.7$, $K = \num{5e9}$; $\Delta G^\circ = -RT\ln K = \qty{-55}{kJ/mol}$. The six Ni--N bonds are of
the same kind on both sides; the number of free molecules rises from four to
seven: the gain of translational entropy drives the reaction.
\end{solution}

\begin{solution}{exo:b3:supramolecular:7}
$\Delta G^\circ = -RT\ln(\num{2.0e5}) = \qty{-30.2}{kJ/mol}$; $T\Delta S^\circ = \Delta H^\circ - \Delta G^\circ = -40 + 30.2 = \qty{-9.8}{kJ/mol}$;
$\Delta S^\circ = \qty{-33}{J.K^{-1}.mol^{-1}}$: enthalpy-driven binding, partly paid back in entropy.
\end{solution}

\begin{solution}{exo:b3:supramolecular:8}
Copper(I) binds two bidentate phenanthroline units and, being tetrahedral, holds
them at right angles, one threaded through the space where the other's ring will
close. Closing both chains into rings then interlocks them. A large excess of
cyanide, which binds copper(I) much more strongly, removes the metal and leaves
the free \omterm{def:b3:supramolecular:interlocked}{catenane}.
\end{solution}

\begin{solution}{exo:b3:supramolecular:9}
$k_c = \pi\Delta\nu/\sqrt2 = \pi \times 120/1.414 = \qty{267}{s^{-1}}$. Eyring:
$\Delta G^\ddagger = RT\ln(k_BT/hk_c) = 8.314 \times 300 \times \ln(\num{6.25e12}/267) = \qty{59.6}{kJ/mol}$.
\end{solution}

\begin{solution}{exo:b3:supramolecular:10}
With excess solid drug, the free drug is fixed at $S_0$. Then $[\mathrm{DC}] = KS_0[\mathrm C]$ and
$C_t = [\mathrm C] + [\mathrm{DC}] = [\mathrm C](1 + KS_0)$, so $[\mathrm{DC}] = KS_0C_t/(1 + KS_0)$ and the total solubility is $S_0 + [\mathrm{DC}]$. Here
$KS_0 = 0.050$: $S = 0.10 + 0.050 \times 10/1.050 = \qty{0.58}{mmol/L}$, nearly six times the intrinsic value.
\end{solution}

\begin{solution}{exo:b3:supramolecular:11}
When $1/K \ll H_0, G_0$, $[\mathrm{HG}] \approx \bigl(H_0 + G_0 - |H_0 - G_0|\bigr)/2 = \min(H_0, G_0)$: the bound fraction is $G_0/H_0$ up to one
equivalent, then 1. The curve no longer depends on $K$: any value large enough
gives the same sharp break within the noise, so the data only show that $K$ exceeds
some minimum.
\end{solution}

\begin{solution}{exo:b3:supramolecular:12}
$k_{\mathrm{off}} = k_{\mathrm{on}}/K = \qty{1e9}{L.mol^{-1}.s^{-1}}/\qty{1500}{L.mol^{-1}} = \qty{6.7e5}{s^{-1}}$. The shift difference
$\qty{0.160}{ppm}$ is \qty{64}{Hz} at \qty{400}{MHz}; coalescence would need $k \approx \pi \times 64/\sqrt2 = \qty{140}{s^{-1}}$. The exchange is
thousands of times faster: one averaged signal.
\end{solution}

\begin{solution}{pb:b3:supramolecular:1}
\textbf{1.} A ring of six $\ce{-CH2CH2O-}$ units; in the complex $\ce{K+}$ sits at the centre, bound by the
six oxygens pointing inwards, the $\ce{CH2}$ groups outside.

\textbf{2.} $\ce{K+}$: $-RT\ln 10 \times 6.1 = \qty{-34.8}{kJ/mol}$; $\ce{Na+}$: \qty{-25.1}{kJ/mol}.

\textbf{3.} $10^{6.1 - 4.4} = 50$.

\textbf{4.} The larger $\ce{K+}$ (\qty{138}{pm}) reaches all six oxygens of the ring at once; the
smaller $\ce{Na+}$ (\qty{102}{pm}) cannot, unless the ring folds, which costs strain.

\textbf{5.} Its outside is a hydrocarbon surface and its charge is buried; $\ce{KCl}$
would need its ions solvated, which toluene cannot do.

\textbf{6.} Water solvates $\ce{K+}$ much better than methanol, and the ion must lose that
water to enter the ring.

\textbf{7.} $\ce{K+}(\text{aq}) + \ce{MnO4-}(\text{aq}) + \mathrm L(\text{org}) \rightleftharpoons \ce{[KL]+MnO4-}(\text{org})$.

\textbf{8.} $[\ce{K+}]_{\mathrm{aq}}$ stays \qty{0.10}{mol/L}: $y = \num{1.0e5} \times 0.10 \times (\num{1.0e-3} - y)^2 = \num{1.0e4}(\num{1.0e-3} - y)^2$.

\textbf{9.} With $u = \num{1.0e-3} - y$: $\num{1.0e4}u^2 + u - \num{1.0e-3} = 0$, $u = \qty{2.70e-4}{mol/L}$, $y = \qty{7.30e-4}{mol/L}$: 73\,\%
extracted.

\textbf{10.} $y = \num{1.0e4}(\num{1.0e-3} - y)(\num{1.0e-2} - y)$ gives $y = \qty{9.89e-4}{mol/L}$: 99\,\%.

\textbf{11.} $y = \num{1.0e4}(\num{1.0e-3} - y)(\num{1.0e-4} - y)$ gives $y = \qty{9.0e-5}{mol/L}$: 9\,\%, limited by the crown,
which is 90\,\% loaded.

\textbf{12.} The permanganate ion keeps its purple colour in the toluene. There
it is unsolvated, a naked anion, and reacts much faster than in water, where a
hydration shell shields it.

\textbf{13.} The sodium ion comes and goes much faster than the difference of the
two resonance frequencies: fast exchange.

\textbf{14.} $\delta = \delta_{\mathrm{free}} + (\delta_{\mathrm{bound}} - \delta_{\mathrm{free}})[\mathrm{HG}]/H_0$, with $[\mathrm{HG}]$ from the exact isotherm, $H_0 = \qty{2.0e-3}{mol/L}$,
$G_0 = (\text{equiv.})\,H_0$, $\delta_{\mathrm{free}} = \qty{3.560}{ppm}$.

\textbf{15.} $K = (1.55 \pm 0.08) \times 10^3\ \unit{L.mol^{-1}}$, $\delta_{\mathrm{bound}} = \qty{3.719 \pm 0.001}{ppm}$.

\textbf{16.} The root-mean-square residual is \qty{0.0014}{ppm}, the size of the reading
precision, with no trend: the 1:1 model fits.

\textbf{17.} Bound fraction $f$ at $G_0 = fH_0 + f/(K(1 - f))$: 0.28 equivalent for $f = 0.2$ and 2.1 for $f = 0.8$.

\textbf{18.} $KH_0 = 10^{6.1} \times \num{2.0e-3} \approx 2500$: the curve would break sharply at one equivalent
and give only a lower bound; a competition experiment or calorimetry at
lower concentration is needed.

\textbf{19.} In the toluene, a homogeneous solution of alkene and permanganate.

\textbf{20.} The reduced manganese leaves the organic phase (as manganese dioxide),
and the crown, freed at the interface, picks up another $\ce{K+}$ and $\ce{MnO4-}$. Each crown
molecule carries permanganate many times: a catalytic amount suffices.

\textbf{21.} $K_{\mathrm{ex}}$ multiplies $[\ce{K+}]$: a high potassium concentration pushes the extraction.

\textbf{22.} Quaternary ammonium salts, such as tetrabutylammonium or
benzyltriethylammonium chloride, whose cation is lipophilic by itself.

\textbf{23.} Benzene causes leukaemia; toluene does the same job with a much smaller
hazard.

\textbf{24.} With \qty{1.0}{mmol/L} of crown, 73\,\% of the permanganate is extracted into the
toluene.
\end{solution}
